\documentclass[11pt]{article}

\usepackage[a4paper,margin=1in]{geometry}
\usepackage{amsmath,amssymb,amsthm,mathtools}
\usepackage{enumitem}
\usepackage{microtype}
\usepackage[dvipsnames]{xcolor}
\usepackage[colorlinks=true,linkcolor=NavyBlue,citecolor=NavyBlue,urlcolor=NavyBlue]{hyperref}

% ---- math shortcuts ----
\newcommand{\E}{\mathbb{E}}
\renewcommand{\Pr}{\mathbb{P}}
\newcommand{\Var}{\operatorname{Var}}
\newcommand{\Z}{\mathbb{Z}}
\newcommand{\given}{\,\vert\,}
\newcommand{\indic}[1]{\mathbf{1}\!\left[#1\right]}
\newcommand{\Bern}{\operatorname{Bernoulli}}
\newcommand{\Bin}{\operatorname{Bin}}

\theoremstyle{definition}
\newtheorem*{problem}{Problem}
\newtheorem*{remark}{Remark}
\newenvironment{solution}{\begin{proof}[Solution]}{\end{proof}}

% source line under each problem title
\newcommand{\source}[3]{\noindent{\small\textit{Source:} #1
  (\href{#2}{QOJ}). \textit{Topics:} #3.}\par}

\setlist[enumerate]{label=(\alph*),leftmargin=*}

\title{\textbf{Discrete Probability Problems}\\[2pt]
\large Nine problems from ICPC-style contests, with solutions}
\author{}
\date{}

\begin{document}
\maketitle

\noindent
These problems come from programming contests. Each statement keeps the
probability core of the original and drops the programming layer (large
inputs, modular output, data structures). Only discrete probability is
used. Problems~1--5 are the main set, in course order; Problems~6--9 are
backups. The selection is explained in \texttt{prob\_course\_candidates.md}.
Every numerical answer below was checked by exact enumeration, and each model
was checked against the original problem's official sample outputs.

\medskip
\noindent
Notation: $H_m=\sum_{i=1}^m \frac1i$ and $H_m^{(2)}=\sum_{i=1}^m \frac1{i^2}$;
$[z^N]\,f(z)$ is the coefficient of $z^N$ in $f(z)$.

\tableofcontents
\clearpage

% =====================================================================
\part*{Main problems}
\addcontentsline{toc}{part}{Main problems}

% ---------------------------------------------------------------------
\section{Mingle}
\source{NAQ 2025, Problem~G}{https://qoj.ac/contest/2554/problem/14443}
  {indicator variables, linearity of expectation, Poisson limit}

\begin{problem}
Let $k\ge 1$ and $n\ge 2k+1$. There are $n$ players and $n$ rooms, both
labelled by $\Z_n=\{0,1,\dots,n-1\}$. The rooms stand in a circle, so room
$n-1$ is next to room $0$. Player $i$ walks into a room chosen uniformly at
random from the $2k+1$ rooms $i-k,\dots,i+k$ (indices mod $n$), independently
of the other players. A player who ends up alone in a room is a
\emph{winner}. Let $W$ be the number of winners.
\begin{enumerate}
  \item Find $\E[W]$. Evaluate it for $n=3$, $k=1$.
  \item Find $\lim_{k\to\infty}\E[W]/n$ and explain the limit with the Poisson
        approximation to the binomial distribution.
\end{enumerate}
\end{problem}

\begin{solution}
(a) Count rooms, not players. A winner is alone in a room, and a room with
exactly one player holds exactly one winner. Hence
\[
  W=\sum_{r\in\Z_n} X_r,\qquad X_r=\indic{\text{room $r$ holds exactly one player}}.
\]
Fix a room $r$. Player $i$ can enter room $r$ only if
$r\in\{i-k,\dots,i+k\}$, that is, $i\in\{r-k,\dots,r+k\}$. These are $2k+1$
distinct players because $n\ge 2k+1$. Each of them enters room $r$ with
probability $\frac{1}{2k+1}$, independently of the others. So the number of
players in room $r$ is $N_r\sim\Bin\bigl(2k+1,\frac1{2k+1}\bigr)$, and
\[
  \Pr(X_r=1)=\Pr(N_r=1)=(2k+1)\cdot\frac{1}{2k+1}\Bigl(1-\frac{1}{2k+1}\Bigr)^{2k}
  =\Bigl(\frac{2k}{2k+1}\Bigr)^{2k}.
\]
The indicators $X_r$ are dependent, but linearity of expectation does not
need independence:
\[
  \E[W]=n\Bigl(\frac{2k}{2k+1}\Bigr)^{2k}.
\]
For $n=3$, $k=1$ this is $3\cdot(2/3)^2=4/3$.

(b) Put $M=2k+1$. Then $\E[W]/n=(1-1/M)^{M-1}\to e^{-1}$ as $k\to\infty$.
The reason is that $N_r\sim\Bin(M,1/M)$ converges in distribution to a
Poisson$(1)$ variable, and $\Pr(\text{Poisson}(1)=1)=e^{-1}$. So in a large
game about $37\%$ of the rooms hold a winner.
\end{solution}

\begin{remark}
The answer does not depend on $n$. If every player instead chooses among all
$n$ rooms, the same argument gives the classical value $n(1-1/n)^{n-1}$.
\end{remark}

% ---------------------------------------------------------------------
\section{Accidental Arithmetic}
\source{NWERC 2025, Problem~A}{https://qoj.ac/contest/2603/problem/14855}
  {linearity of expectation, symmetry, conditioning}

\begin{problem}
A positive integer $N$ has decimal digits $d_1d_2\cdots d_L$ ($d_1\ne 0$).
You type the digits from left to right into a faulty calculator. Into each of
the $L-1$ gaps between consecutive digits it inserts, independently, a ``$+$''
with probability $\beta$, a ``$-$'' with probability $\gamma$, and nothing with
probability $q=1-\beta-\gamma$, where $\beta+\gamma>0$. The resulting
expression, such as $12-3+45$, is then evaluated; a term may begin with zeros,
so $1+05=6$. Let $V$ be the value, and let $P_l$ be the number formed by the
first $l$ digits (so $P_L=N$).
\begin{enumerate}
  \item Let $\beta=\gamma=\alpha$. Show that
  \[
    \E[V]=\sum_{l=1}^{L-1}2\alpha(1-2\alpha)^{l-1}P_l+(1-2\alpha)^{L-1}N,
  \]
  and evaluate it for $N=12345$, $\alpha=0.45$.
  \item Now let $\beta,\gamma$ be arbitrary. Let $F$ be the first term of the
  expression and $S$ the sum of all its terms, each taken with a plus sign.
  Show that $\E[V]=(1-\mu)\E[F]+\mu\,\E[S]$ with
  $\mu=\frac{\beta-\gamma}{\beta+\gamma}$, and compute $\E[F]$ and $\E[S]$.
\end{enumerate}
\end{problem}

\begin{solution}
Number the gaps $1,\dots,L-1$, gap $g$ lying between $d_g$ and $d_{g+1}$.
Call a gap a \emph{sign gap} if it receives a sign, and let $A$ be the
(random) set of sign gaps. The set $A$ cuts the digits into terms
$F=T_0,T_1,\dots,T_s$, and
\[
  V=F+\sum_{t=1}^{s}\sigma_tT_t,
\]
where $\sigma_t\in\{+1,-1\}$ is the sign in front of $T_t$. Given $A$, the
terms $T_t$ are fixed numbers, and the signs at the gaps in $A$ are independent,
each equal to $+1$ with probability $\beta/(\beta+\gamma)$. So
$\E[\sigma_t\given A]=\mu$ and
\[
  \E[V\given A]=F+\mu\sum_{t=1}^{s}T_t=(1-\mu)F+\mu S.
\]
Taking expectations gives $\E[V]=(1-\mu)\E[F]+\mu\E[S]$, which is the identity
in (b).

\emph{The first term.} Let $G$ be the first sign gap, with $G=L$ if there is
none. Then $F=P_G$, $\Pr(G=l)=(1-q)q^{l-1}$ for $l<L$, and
$\Pr(G=L)=q^{L-1}$. Hence
\[
  \E[F]=\sum_{l=1}^{L-1}(1-q)q^{l-1}P_l+q^{L-1}N.
\]

(a) If $\beta=\gamma=\alpha$ then $\mu=0$: every later term is equally likely
to be added or subtracted, whatever its digits, so only the first term
survives. With $1-q=2\alpha$ the formula above is the claimed one. For
$N=12345$ and $\alpha=0.45$ we have $q=0.1$, and
\begin{align*}
  \E[V]&=0.9\cdot 1+0.09\cdot 12+0.009\cdot 123+0.0009\cdot 1234+0.0001\cdot 12345\\
  &=0.9+1.08+1.107+1.1106+1.2345=5.4321.
\end{align*}

(b) It remains to find $\E[S]$; use linearity over digits. Digit $d_l$
contributes $d_l\cdot 10^{J_l}$ to $S$, where $J_l$ is the number of digits
after position $l$ in the same term. That is, $J_l$ counts the gaps
$l,l+1,\dots$ that are empty before the first sign gap, capped at $L-l$. So
$\Pr(J_l=j)=(1-q)q^j$ for $0\le j<L-l$ and $\Pr(J_l=L-l)=q^{L-l}$, and
\[
  \E[S]=\sum_{l=1}^{L}d_l\Bigl(\sum_{j=0}^{L-l-1}(1-q)(10q)^j+(10q)^{L-l}\Bigr).
\]
Check: if $q=0$, every gap holds a sign, so $\E[F]=d_1$, $\E[S]=\sum_l d_l$
and $\E[V]=d_1+\mu\sum_{l\ge 2}d_l$, as it should be.
\end{solution}

% ---------------------------------------------------------------------
\section{502 Bad Gateway}
\source{ICPC Asia EC Online 2024 (II), Problem~L}{https://qoj.ac/contest/1799/problem/9381}
  {geometric distribution, optimal stopping}

\begin{problem}
Fix an integer $T\ge 1$. A timer is set to a value drawn uniformly from
$\{1,\dots,T\}$ and counts down by $1$ each second; you are done when it
reaches $0$. At the end of every second at which the timer is still positive,
you may \emph{reset} it: it is then set to a fresh uniform value from
$\{1,\dots,T\}$, independent of the past. You always see the timer. You want
to minimise the expected time until you are done.
\begin{enumerate}
  \item For $1\le c\le T$, consider the policy ``keep a drawn value if it is at
  most $c$; otherwise wait one second and reset''. Show that its expected time
  is
  \[
    g(c)=\frac{T}{c}+\frac{c-1}{2}.
  \]
  \item Show that some policy of this form is optimal.
  \item Find the optimal $c$ and the minimal expected time. Evaluate it for
  $T=3$, $10$ and $1000$, and find its asymptotics as $T\to\infty$.
\end{enumerate}
\end{problem}

\begin{solution}
(a) Each draw is kept with probability $c/T$, independently of the other
draws, so the number of rejected draws is geometric with mean
$\frac{1-c/T}{c/T}=\frac{T-c}{c}$. A rejected draw is larger than $c\ge 1$, so
the timer is still positive after one second, and the draw costs exactly one
second. The kept draw is uniform on $\{1,\dots,c\}$ and costs its value, with
mean $\frac{c+1}{2}$. Hence
\[
  g(c)=\frac{T-c}{c}+\frac{c+1}{2}=\frac{T}{c}+\frac{c-1}{2}.
\]

(b) Let $V$ be the minimal expected time measured from the moment a fresh
value is drawn; $V\le g(T)=\frac{T+1}{2}$ is finite. After a reset the future
is a fresh copy of the whole problem, so the best you can do from a reset is
$V$, whatever happened before.

Suppose that at the end of some second the timer shows $s\ge 1$. If you never
reset again, you need $s$ more seconds. If you reset after waiting $w\ge 0$
more seconds, you need at least $w+V\ge V$, with equality for $w=0$. So the
best you can do from this point is $\min(s,V)$: reset now if $s>V$, and never
reset if $s\le V$.

After a draw $u\ge 2$, the first decision point shows $s=u-1$. So it is
optimal to keep $u$ exactly when $u\le V+1$. Once a value is kept, later
decision points show smaller values, which are still at most $V$, so it stays
kept. (A draw $u=1$ ends the game at once.) This is the threshold policy with
$c=\min\bigl(T,\lfloor V+1\rfloor\bigr)$. Therefore $V=\min_{1\le c\le T}g(c)$.

(c) For $1\le c<T$,
\[
  g(c+1)-g(c)=\frac12-\frac{T}{c(c+1)}.
\]
So $g$ decreases while $c(c+1)<2T$ and does not decrease afterwards. The
optimal threshold is the least $c$ with $c(c+1)\ge 2T$, namely
$c^*=\bigl\lceil(\sqrt{8T+1}-1)/2\bigr\rceil$; when $c^*(c^*+1)=2T$ the values
$c^*$ and $c^*+1$ tie.
\begin{itemize}
  \item $T=3$: $c^*=2$ (tied with $3$), $V=\frac32+\frac12=2$.
  \item $T=10$: $c^*=4$ (tied with $5$), $V=\frac{10}{4}+\frac32=4$.
  \item $T=1000$: $44\cdot 45=1980<2000\le 45\cdot 46$, so $c^*=45$ and
        $V=\frac{1000}{45}+22=\frac{398}{9}\approx 44.22$.
\end{itemize}
For large $T$, the function $x\mapsto T/x+x/2$ is minimised at $x=\sqrt{2T}$
with value $\sqrt{2T}$, and rounding $x$ to an integer changes the value by
$O(T^{-1/2})$. Hence $c^*=\sqrt{2T}+O(1)$ and
$V=\sqrt{2T}-\frac12+O(T^{-1/2})$; for $T=1000$ this gives $44.22$. Never
resetting would cost $\frac{T+1}{2}$, so resets are worth much more than the
second each one wastes.
\end{solution}

% ---------------------------------------------------------------------
\section{Uniform Chemistry}
\source{SWERC 2022--2023, Problem~B}{https://qoj.ac/contest/1123/problem/5646}
  {Markov chains, independent indicators, random permutations}

\begin{problem}
Fix integers $1\le a<n$ and put $m=n-a$. A particle starts at $a$. At each
step, if it is at $x<n$, it jumps to a point chosen uniformly from
$\{x+1,\dots,n\}$, independently of the past. Let $T$ be the number of steps
until it reaches $n$.
\begin{enumerate}
  \item For $a<j<n$, let $I_j$ be the indicator that the particle visits $j$.
  Show that $\Pr(I_j=1)=\frac{1}{n-j+1}$ and that $I_{a+1},\dots,I_{n-1}$ are
  independent.
  \item Find $\E[T]$ and $\Var(T)$.
  \item Let $c(m,t)$ be the unsigned Stirling numbers of the first kind,
  defined by $\sum_t c(m,t)z^t=z(z+1)\cdots(z+m-1)$; $c(m,t)$ is the number of
  permutations of $m$ elements with exactly $t$ cycles. Show that
  $\Pr(T=t)=c(m,t)/m!$.
  \item Three researchers each run an independent copy of this process with
  $a=1$ and $n=3$, one step per year. A researcher wins if nobody reaches $n$
  strictly earlier, so ties all win. Find the probability that a given
  researcher wins.
\end{enumerate}
\end{problem}

\begin{solution}
(a) Compute the joint law directly. Let
$S=\{s_1<\dots<s_r\}\subseteq\{a+1,\dots,n-1\}$. The particle visits exactly
the points of $S$ between $a$ and $n$ if and only if its path is
$a\to s_1\to\dots\to s_r\to n$, which has probability
\[
  \frac{1}{n-a}\prod_{i=1}^{r}\frac{1}{n-s_i}.
\]
Now put $q_j=\frac{1}{n-j+1}$. The product
$\prod_{j=a+1}^{n-1}(1-q_j)=\prod_{j=a+1}^{n-1}\frac{n-j}{n-j+1}$ telescopes to
$\frac{1}{n-a}$, and $\frac{q_j}{1-q_j}=\frac{1}{n-j}$. Hence
\[
  \prod_{j\in S}q_j\prod_{j\notin S}(1-q_j)
  =\prod_{j=a+1}^{n-1}(1-q_j)\cdot\prod_{j\in S}\frac{q_j}{1-q_j}
  =\frac{1}{n-a}\prod_{j\in S}\frac{1}{n-j},
\]
which is the same probability. So $(I_{a+1},\dots,I_{n-1})$ has the law of
independent $\Bern(q_j)$ variables, and both claims follow.

There is also a direct argument for the marginal: the particle must enter
$\{j,\dots,n\}$ at some step, and the jump that first enters this set lands
uniformly on its $n-j+1$ points.

(b) Each step lands on a new point and the last step lands on $n$, so
$T=1+\sum_{j=a+1}^{n-1}I_j$. Substitute $i=n-j+1$, which runs over
$2,\dots,m$:
\[
  T=1+\sum_{i=2}^{m}B_i,\qquad B_i\sim\Bern(1/i)\ \text{independent}.
\]
Hence $\E[T]=1+\sum_{i=2}^m\frac1i=H_m$ and, by independence,
\[
  \Var(T)=\sum_{i=2}^m\frac1i\Bigl(1-\frac1i\Bigr)=H_m-H_m^{(2)}
\]
(the term $i=1$ of $H_m-H_m^{(2)}$ is zero).

(c) By independence,
\[
  \E\bigl[z^T\bigr]=z\prod_{i=2}^{m}\Bigl(1-\frac1i+\frac zi\Bigr)
  =z\prod_{i=2}^{m}\frac{z+i-1}{i}=\frac{z(z+1)\cdots(z+m-1)}{m!}
  =\sum_t\frac{c(m,t)}{m!}z^t.
\]
Comparing coefficients gives $\Pr(T=t)=c(m,t)/m!$. So $T$ has the law of the
number of cycles of a uniformly random permutation of $m$ elements.

(d) With $a=1$ and $n=3$, the first jump goes to $3$ with probability
$\frac12$ (then $T=1$); otherwise it goes to $2$ and then to $3$ (then $T=2$).
A given researcher wins if their $T=1$, or if their $T=2$ and both other
researchers also have $T=2$:
\[
  \Pr(\text{win})=\frac12+\Bigl(\frac12\Bigr)^3=\frac58.
\]
In general, researcher $i$ wins with probability
$\sum_t\Pr(T_i=t)\prod_{j\ne i}\Pr(T_j\ge t)$, and (c) gives each
distribution.
\end{solution}

% ---------------------------------------------------------------------
\section{Jinx or Jackpot}
\source{NERC 2025, Problem~J}{https://qoj.ac/contest/2785/problem/15902}
  {Bayes' rule, dynamic programming, value of information}

\begin{problem}
A slot machine pays even odds: if you bet $x$, you win $x$ with probability
$p$ (a \emph{jackpot}) and lose $x$ with probability $1-p$ (a \emph{jinx}).
Given $p$, the rounds are independent. You start with \$1000 and may play at
most $k$ rounds. In each round you choose a bet $x$ with
$0\le x\le{}$(current balance); a bet of $0$ is allowed, and you still see
whether that round was a jackpot. You want to maximise the expected final
balance.
\begin{enumerate}
  \item Show that the optimal expected final balance from balance $w$ is $w$
  times a number that depends only on the outcomes seen so far, and that in
  every round an optimal bet is either $0$ or the whole balance.
  \item Suppose $p$ is known. Find an optimal strategy and the maximal expected
  profit. Evaluate the profit for $p=0.61$, $k=5$.
  \item Suppose $p$ is unknown: it is $0.7$ or $0.3$, each with prior
  probability $\frac12$. Let $k=2$. Find an optimal strategy and the maximal
  expected profit, and compare it with the best strategy that bets $0$ in the
  first round.
\end{enumerate}
\end{problem}

\begin{solution}
(a) Let $h$ be the sequence of outcomes seen so far, and let $j$ be the number
of rounds left. Given $h$, the next round is a jackpot with probability
$\hat p(h)=\E[p\given h]$, the posterior mean of $p$ (law of total
probability). Write $hW$ and $hL$ for $h$ followed by a jackpot or a jinx.
Let $V_j(w,h)$ be the optimal expected final balance. We show by induction on
$j$ that $V_j(w,h)=w\,v_j(h)$, where $v_0\equiv 1$ and
\[
  v_j(h)=\max\Bigl\{\hat p\,v_{j-1}(hW)+(1-\hat p)\,v_{j-1}(hL),\;
  2\hat p\,v_{j-1}(hW)\Bigr\},\qquad \hat p=\hat p(h).
\]
For $j=0$ this is clear. For $j\ge 1$, a bet $x\in[0,w]$ is followed by
balance $w+x$ after a jackpot and $w-x$ after a jinx, and the outcome is seen
either way, so by induction
\[
  V_j(w,h)=\max_{0\le x\le w}\Bigl[\hat p\,(w+x)\,v_{j-1}(hW)
  +(1-\hat p)\,(w-x)\,v_{j-1}(hL)\Bigr].
\]
The bracket is linear in $x$, so its maximum is attained at $x=0$ or $x=w$.
These two choices give the two expressions in the formula for $v_j$.

(b) If $p$ is known, then $\hat p=p$ in every round and
$v_j=v_{j-1}\max(1,2p)$, so $v_k=\max(1,2p)^k$. If $p>\frac12$, go all-in
every round; the expected profit is $1000\bigl((2p)^k-1\bigr)$. If
$p\le\frac12$, never bet; the profit is $0$. For $p=0.61$ and $k=5$ the
profit is $1000\,(1.22^5-1)=1000\cdot 1.7027081632\approx\$1702.71$.

(c) Bayes' rule after one jackpot:
\[
  \Pr(p=0.7\given W)=\frac{\frac12\cdot 0.7}{\frac12\cdot 0.7+\frac12\cdot 0.3}=0.7,
  \qquad \hat p(W)=0.7\cdot 0.7+0.3\cdot 0.3=0.58.
\]
Likewise $\Pr(p=0.7\given L)=0.3$ and $\hat p(L)=0.3\cdot 0.7+0.7\cdot
0.3=0.42$. Before any round, $\hat p=0.5$.

\emph{Round 2.} After a jackpot, $v_1(W)=\max(1,\,2\cdot 0.58)=1.16$: go
all-in. After a jinx, $v_1(L)=\max(1,\,2\cdot 0.42)=1$: bet nothing.

\emph{Round 1.} Betting $0$ gives $\frac12\cdot 1.16+\frac12\cdot 1=1.08$.
Going all-in gives $2\cdot\frac12\cdot 1.16=1.16$. So the optimal strategy is
to go all-in in round 1 and, after a jackpot, all-in again. The maximal
expected profit is $1000\,(1.16-1)=\$160$. The best strategy that only watches
round 1 earns $1000\,(1.08-1)=\$80$.

Round 1 is a fair bet ($\hat p=\frac12$), yet staking everything on it is
optimal. A jackpot both doubles the balance and makes the machine look better,
so the balance is largest exactly when the next bet is most favourable.
\end{solution}

% =====================================================================
\part*{Backup problems}
\addcontentsline{toc}{part}{Backup problems}

% ---------------------------------------------------------------------
\section{Orchard Fence}
\source{2025 ICPC Southeast USA Regional, Division~1, Problem~A; listed under
Mid-Atlantic USA in \texttt{data/tagged.json}}{https://qoj.ac/contest/3502/problem/16814}
  {linearity of expectation over edges, counting}

\begin{problem}
Let $z_1,\dots,z_n$ be $n\ge 3$ points in the plane in general position (no
three on a line). A uniformly random subset $K$ of $k$ of them is kept, where
$2\le k\le n$. Let $L$ be the perimeter of the convex hull of $K$; for $k=2$
the hull is a segment, and its perimeter is twice its length.
\begin{enumerate}
  \item For an ordered pair $(i,j)$ with $i\ne j$, let $\ell_{ij}$ be the number
  of the other $n-2$ points that lie strictly to the left of the directed line
  from $z_i$ to $z_j$. Show that
  \[
    \E[L]=\frac{1}{\binom nk}\sum_{i\ne j}|z_iz_j|\binom{\ell_{ij}}{k-2}.
  \]
  \item Evaluate $\E[L]$ for the six vertices of a regular hexagon with side
  $1$, for $k=3$ and for $k=4$.
\end{enumerate}
\end{problem}

\begin{solution}
(a) Traverse the hull boundary counterclockwise. Then
\[
  L=\sum_{i\ne j}|z_iz_j|\cdot\indic{z_i\to z_j\text{ is a counterclockwise hull edge of }K}.
\]
For $k\ge 3$ in general position, $z_i\to z_j$ is a counterclockwise hull edge
exactly when $i,j\in K$ and every other point of $K$ lies strictly to the left
of the line $z_i\to z_j$. The number of such $k$-subsets is
$\binom{\ell_{ij}}{k-2}$: choose the other $k-2$ kept points among the
$\ell_{ij}$ points on the left. For $k=2$ the formula counts the segment in
both directions, since $\binom{\ell}{0}=1$, which matches the convention.
Linearity of expectation gives the formula.

(b) Label the vertices $0,\dots,5$ counterclockwise. There are three kinds of
pairs:
\begin{itemize}
  \item 6 adjacent pairs, length $1$: $\ell=4$ in the counterclockwise
  direction and $\ell=0$ in the other;
  \item 6 pairs one vertex apart, length $\sqrt3$: $\ell=3$ and $\ell=1$;
  \item 3 opposite pairs, length $2$: $\ell=2$ in both directions.
\end{itemize}
For $k=3$, the weights $\binom{\ell}{1}$ of the two directions add up to $4$
for every pair, so
\[
  \E[L]=\frac{4\,(6\cdot 1+6\sqrt3+3\cdot 2)}{\binom63}=\frac{48+24\sqrt3}{20}
  =\frac{12+6\sqrt3}{5}\approx 4.478.
\]
Check: a triangle's perimeter is the sum of its three sides, so $\E[L]$ is
three times the mean distance of a random pair, $3\cdot(12+6\sqrt3)/15$.

For $k=4$, the weights $\binom{\ell}{2}$ add up to $6+0$, $3+0$ and $1+1$ for
the three kinds, so
\[
  \E[L]=\frac{6\cdot 1\cdot 6+6\sqrt3\cdot 3+3\cdot 2\cdot 2}{\binom64}
  =\frac{48+18\sqrt3}{15}=\frac{16+6\sqrt3}{5}\approx 5.278.
\]
\end{solution}

\begin{remark}
In the original problem the kept points are tree trunks, disks of diameter
$d$. The shortest fence around them has length $L+\pi d$, so the expected
fence length is $\E[L]+\pi d$.
\end{remark}

% ---------------------------------------------------------------------
\section{Chopsticks}
\source{2025 ICPC Asia Taiwan Online Programming Contest, Problem~H}{https://qoj.ac/contest/2621/problem/15121}
  {parity, indicators, generating functions}

\begin{problem}
A box holds chopsticks of $m$ types, $k_i$ of type $i$, and $s=\sum_i k_i$ in
total. You draw $2n\le s$ of them uniformly at random without replacement and
hand them out in $n$ pairs so that as many pairs as possible are of one type.
Let $M$ be the number of mixed pairs, and let $X_i$ be the number of drawn
chopsticks of type $i$.
\begin{enumerate}
  \item Show that $M=\frac12\,\#\{i:X_i\text{ odd}\}$, and hence
  $\E[M]=\frac12\sum_i\Pr(X_i\text{ odd})$.
  \item Show that $\Pr(X_i\text{ odd})=\frac12\bigl(1-\E[(-1)^{X_i}]\bigr)$ and
  \[
    \E\bigl[(-1)^{X_i}\bigr]=\frac{[z^{2n}]\,(1-z)^{k_i}(1+z)^{s-k_i}}{\binom{s}{2n}}.
  \]
  \item Find $\E[M]$ for $8$ red and $8$ blue chopsticks and $n=5$.
  \item Suppose instead that the $2n$ chopsticks are drawn independently, each
  of type $i$ with probability $p_i$. Find $\E[M]$.
\end{enumerate}
\end{problem}

\begin{solution}
(a) Type $i$ gives at most $\lfloor X_i/2\rfloor$ same-type pairs, and pairing
within types achieves $\sum_i\lfloor X_i/2\rfloor$. The leftover chopsticks,
one of each type with $X_i$ odd, are then paired arbitrarily. If $O$ is the
number of odd $X_i$, the matched pairs number $\sum_i\lfloor X_i/2\rfloor=(2n-O)/2$,
so $M=n-(2n-O)/2=O/2$. Linearity gives
$\E[M]=\frac12\sum_i\Pr(X_i\text{ odd})$.

(b) Since $(-1)^X=1-2\cdot\indic{X\text{ odd}}$, we get
$\Pr(X\text{ odd})=\frac12\bigl(1-\E[(-1)^X]\bigr)$. The count $X_i$ is
hypergeometric:
$\Pr(X_i=x)=\binom{k_i}{x}\binom{s-k_i}{2n-x}\big/\binom{s}{2n}$. So
\[
  \binom{s}{2n}\E\bigl[(-1)^{X_i}\bigr]=\sum_x(-1)^x\binom{k_i}{x}\binom{s-k_i}{2n-x}
  =[z^{2n}]\,(1-z)^{k_i}(1+z)^{s-k_i},
\]
because multiplying the two generating functions
$(1-z)^{k_i}=\sum_x(-1)^x\binom{k_i}{x}z^x$ and
$(1+z)^{s-k_i}=\sum_y\binom{s-k_i}{y}z^y$ collects exactly these products in
the coefficient of $z^{2n}$.

(c) Here $k=8$, $s=16$ and $2n=10$. Since $(1-z)^8(1+z)^8=(1-z^2)^8$,
\[
  [z^{10}]\,(1-z^2)^8=\binom85(-1)^5=-56,\qquad
  \E\bigl[(-1)^{X_{\text{red}}}\bigr]=\frac{-56}{\binom{16}{10}}=\frac{-56}{8008}=-\frac{1}{143}.
\]
So $\Pr(X_{\text{red}}\text{ odd})=\frac12\bigl(1+\frac1{143}\bigr)=\frac{72}{143}$,
and the same holds for blue. Hence $\E[M]=\frac12\cdot 2\cdot\frac{72}{143}
=\frac{72}{143}\approx 0.503$. (With two colours, $M$ is simply the indicator
that $X_{\text{red}}$ is odd.)

(d) Now $X_i\sim\Bin(2n,p_i)$, so
$\E[(-1)^{X_i}]=\sum_x\binom{2n}{x}(-p_i)^x(1-p_i)^{2n-x}=(1-2p_i)^{2n}$, and
\[
  \E[M]=\frac14\sum_i\Bigl(1-(1-2p_i)^{2n}\Bigr).
\]
\end{solution}

% ---------------------------------------------------------------------
\section{Crashing Competition Computer}
\source{NWERC 2022, Problem~C}{https://qoj.ac/contest/1044/problem/5127}
  {consecutive successes, recursions for expectations, convexity}

\begin{problem}
You must type a text of $c$ characters, one time unit per character. After
each typed character, independently, the computer crashes with probability
$p\in(0,1)$. A crash costs $r$ time units of recovery, after which you continue
from the last save; the start of the text counts as saved. A save takes $t$
time units and never crashes, and the text must be saved after its last
character. Put $q=1-p$.
\begin{enumerate}
  \item Show that typing a block of $k$ characters that starts right after a
  save and ends with a save takes expected time $h(k)+t$, where
  \[
    h(k)=\Bigl(\frac1p+r\Bigr)\bigl(q^{-k}-1\bigr).
  \]
  \item Show that, for a fixed number of saves, the expected total time is
  smallest when the block lengths differ by at most one.
  \item Find the minimal expected time for $c=3$, $p=\frac12$, $r=2$, $t=5$.
\end{enumerate}
\end{problem}

\begin{solution}
(a) Let $E_j$ be the expected time, starting from the save, until $j$
characters of the block have been typed with no crash in between. Once $k-1$
characters are typed, you type one more (one time unit). With probability $p$
it crashes; you then pay $r$ and are back at the save, so you need $E_k$ more
in expectation. Hence
\[
  E_k=E_{k-1}+1+p\,(r+E_k),\qquad\text{so}\qquad E_k=\frac{E_{k-1}+1+pr}{q},
  \qquad E_0=0.
\]
Unrolling the recursion,
\[
  E_k=(1+pr)\sum_{i=1}^{k}q^{-i}=(1+pr)\,\frac{q^{-k}-1}{1-q}
  =\Bigl(\frac1p+r\Bigr)\bigl(q^{-k}-1\bigr)=h(k),
\]
and the final save adds $t$.

(b) First, a fixed plan of save positions is enough. A crash returns you to
the state you were in right after the last save, so where to save next from a
given save position can be decided once and for all. The plan is a sequence
$0=y_0<y_1<\dots<y_B=c$ of save positions, with block lengths
$k_b=y_b-y_{b-1}$. By (a) and linearity its expected time is
$\sum_{b=1}^{B}h(k_b)+Bt$.

The increments $\Delta(k)=h(k+1)-h(k)=\bigl(\frac1p+r\bigr)\frac pq\,q^{-k}$
increase with $k$, so $h$ is convex. If two blocks have lengths $a\ge b+2$,
replacing them by $a-1$ and $b+1$ changes the total by
\[
  \Delta(b)-\Delta(a-1)\le 0,
\]
because $b<a-1$. Repeating this until all lengths differ by at most one
never increases the total. Writing $c=Bu+v$ with $0\le v<B$, the best plan
with $B$ saves therefore costs $v\,h(u+1)+(B-v)\,h(u)+Bt$, and it remains to
minimise over $B\in\{1,\dots,c\}$.

(c) Here $\frac1p+r=4$ and $q^{-1}=2$, so $h(k)=4(2^k-1)$: $h(1)=4$,
$h(2)=12$, $h(3)=28$.
\begin{itemize}
  \item One save (block $3$): $28+5=33$.
  \item Two saves (blocks $2,1$): $12+4+2\cdot 5=26$.
  \item Three saves (blocks $1,1,1$): $3\cdot(4+5)=27$.
\end{itemize}
The minimum is $26$: save once after the second (or the first) character, and
again at the end.
\end{solution}

% ---------------------------------------------------------------------
\section{Disbursement on Quarantine Policy}
\source{ICPC Asia Taiwan 2024, Problem~D}{https://qoj.ac/contest/2622/problem/15129}
  {linearity of expectation, complements, independence}

\begin{problem}
Passengers sit in a rectangular grid of seats, and each is infected
independently with probability $\frac12$. An infected passenger is quarantined
for $d_0$ days. A healthy passenger with at least one infected edge-neighbour
(sharing a side) is quarantined for $d_1$ days. A healthy passenger with no
infected edge-neighbour but at least one infected corner-neighbour (sharing
only a corner) is quarantined for $d_2$ days. Everyone else is not
quarantined.
\begin{enumerate}
  \item A seat has $e$ edge-neighbours and $c$ corner-neighbours. Find the
  expected quarantine length of its passenger.
  \item Find the expected total number of quarantine days for a $3\times 3$
  grid with $d_0=4$, $d_1=2$, $d_2=1$.
\end{enumerate}
\end{problem}

\begin{solution}
(a) The events ``the passenger is infected'', ``no edge-neighbour is
infected'' and ``no corner-neighbour is infected'' concern disjoint sets of
seats, so they are independent, with probabilities $\frac12$, $2^{-e}$ and
$2^{-c}$. Using complements,
\begin{align*}
  \Pr(\text{healthy, some edge-neighbour infected})&=\tfrac12\bigl(1-2^{-e}\bigr),\\
  \Pr(\text{healthy, no edge-neighbour infected, some corner-neighbour infected})
  &=\tfrac12\,2^{-e}\bigl(1-2^{-c}\bigr).
\end{align*}
So the expected quarantine length is
\[
  \frac{d_0}{2}+\frac{d_1}{2}\bigl(1-2^{-e}\bigr)+\frac{d_2}{2}\,2^{-e}\bigl(1-2^{-c}\bigr).
\]

(b) By linearity, add up the seats. With $d_0=4$, $d_1=2$, $d_2=1$:
\begin{itemize}
  \item 4 corner seats ($e=2$, $c=1$): $2+\frac34+\frac1{16}=\frac{45}{16}$;
  \item 4 side seats ($e=3$, $c=2$): $2+\frac78+\frac3{64}=\frac{187}{64}$;
  \item 1 centre seat ($e=4$, $c=4$): $2+\frac{15}{16}+\frac{15}{512}=\frac{1519}{512}$.
\end{itemize}
The expected total is
\[
  4\cdot\frac{45}{16}+4\cdot\frac{187}{64}+\frac{1519}{512}
  =\frac{5760+5984+1519}{512}=\frac{13263}{512}\approx 25.90.
\]
\end{solution}

\begin{remark}
In the original problem some passengers are known to be infected or healthy.
The same argument works with each seat's infection probability in
$\{0,\frac12,1\}$; the probability that no edge-neighbour is infected becomes
the product of $1-\Pr(\text{infected})$ over the edge-neighbours.
\end{remark}

\end{document}
